\section{Teorema fundamental del cálculo}

Para \(f\) integrable en \([a, b]\) y \(c, d \in [a, b]\) con \(c > d\) convenimos \(\int_c^d f(x) \,\dd x = -\int_d^c f(x) \,\dd x\), y \(\int_c^c f(x) \,\dd x = 0\). Con esta convención vale \(\int_a^b f = \int_a^c f + \int_c^b f\) para \(a, b, c\) cualesquiera en el dominio.

\begin{theorem}[Teorema fundamental del cálculo]
  Sea \(f\colon [a, b] \to \R\) integrable. Si \(f\) es continua en \(c \in [a, b]\), entonces \(F\colon [a, b] \to \R\), \(F(x) = \int_a^x f(t) \,\dd t\), es derivable en \(c\) (con derivadas laterales en los extremos) y \(F'(c) = f(c)\).
\end{theorem}

\begin{proof}
  Sea \(\e > 0\). Como \(f\) es continua en \(c\), existe \(\delta > 0\) tal que \(t \in [a, b]\) y \(|t - c| < \delta\) implican \(|f(t) - f(c)| < \e\). Sea \(h \neq 0\) con \(|h| < \delta\) y \(c + h \in [a, b]\). Como \(F(c + h) - F(c) = \int_c^{c+h} f(t) \,\dd t\) y \(\int_c^{c+h} f(c) \,\dd t = h f(c)\),
  \[
    \left| \frac{F(c + h) - F(c)}{h} - f(c) \right| = \frac{1}{|h|} \left| \int_c^{c+h} \bigl( f(t) - f(c) \bigr) \dd t \right| \leq \frac{1}{|h|} \cdot \e |h| = \e.
  \]
  Luego \(\lim_{h \to 0} \frac{F(c + h) - F(c)}{h} = f(c)\).
\end{proof}

\begin{corollary}[Regla de Barrow]
  Sea \(f\colon [a, b] \to \R\) integrable y sea \(F\colon [a, b] \to \R\) derivable con \(F' = f\). Entonces \(\int_a^b f(x) \,\dd x = F(b) - F(a)\).
\end{corollary}

\begin{proof}
  Sea \(P = \{ t_0, \dots, t_n \}\) una partición de \([a, b]\). Por el teorema del valor medio, para cada \(i\) existe \(\alpha_i \in (t_{i-1}, t_i)\) con \(F(t_i) - F(t_{i-1}) = f(\alpha_i)(t_i - t_{i-1})\). Entonces
  \[
    F(b) - F(a) = \sum_{i=1}^{n} \bigl( F(t_i) - F(t_{i-1}) \bigr) = \sum_{i=1}^{n} f(\alpha_i)(t_i - t_{i-1}).
  \]
  Como \(m_i \leq f(\alpha_i) \leq M_i\), resulta \(s(f, P) \leq F(b) - F(a) \leq S(f, P)\) para toda partición \(P\). Luego \(\lowint_a^b f \,\dd x \leq F(b) - F(a) \leq \upint_a^b f \,\dd x\), y como \(f\) es integrable, \(F(b) - F(a) = \int_a^b f(x) \,\dd x\).
\end{proof}

\begin{definition}[Norma de una partición]
  Sea \(P = \{ t_0, \dots, t_n \}\) una partición de \([a, b]\). La \emph{norma} de \(P\) es \(\norm{P} = \max_{1 \leq i \leq n} (t_i - t_{i-1})\).
\end{definition}

\begin{theorem}
  Sea \(f\colon [a, b] \to \R\) acotada. Para todo \(\e > 0\) existe \(\delta > 0\) tal que si \(\norm{P} < \delta\), entonces \(S(f, P) < \upint_a^b f(x) \,\dd x + \e\).
\end{theorem}

\begin{proof}
  Supongamos primero \(f \geq 0\). Sea \(M = \sup f\). Dado \(\e > 0\) existe una partición \(P_0 = \{ t_0, \dots, t_n \}\) con \(S(f, P_0) < \upint_a^b f \,\dd x + \frac{\e}{2}\). Sea \(\delta = \frac{\e}{2n (M + 1)}\) y sea \(P\) una partición con \(\norm{P} < \delta\).

  Dividimos los intervalos de \(P\) en dos tipos: los contenidos en algún \([t_{i-1}, t_i]\) y los restantes. Cada intervalo del segundo tipo contiene algún \(t_i\) en su interior, así que hay a lo sumo \(n\) de ellos, y cada uno tiene longitud menor que \(\delta\) y \(M_k \leq M\). Para cada \(i\), los intervalos del primer tipo contenidos en \([t_{i-1}, t_i]\) tienen supremo \(\leq M_i\) y longitudes que suman a lo sumo \(t_i - t_{i-1}\), de modo que su aporte a \(S(f, P)\) es a lo sumo \(M_i (t_i - t_{i-1})\) (pues \(M_i \geq 0\)). Entonces
  \[
    S(f, P) \leq \sum_{i=1}^{n} M_i (t_i - t_{i-1}) + n M \delta < S(f, P_0) + \frac{\e}{2} < \upint_a^b f \,\dd x + \e.
  \]

  Para \(f\) acotada cualquiera, existe \(c \in \R\) con \(g = f + c \geq 0\). Como \(S(g, P) = S(f, P) + c(b - a)\) para toda partición y \(\upint_a^b g \,\dd x = \upint_a^b f \,\dd x + c(b - a)\), el resultado para \(g\) implica el de \(f\).
\end{proof}

\begin{corollary}
  Si \(f\colon [a, b] \to \R\) es acotada, entonces
  \[
    \lim_{\norm{P} \to 0} S(f, P) = \upint_a^b f(x) \,\dd x, \qquad \lim_{\norm{P} \to 0} s(f, P) = \lowint_a^b f(x) \,\dd x.
  \]
  En particular, si \(f\) es integrable, \(\int_a^b f(x) \,\dd x = \lim_{\norm{P} \to 0} S(f, P) = \lim_{\norm{P} \to 0} s(f, P)\).
\end{corollary}

\begin{proof}
  La primera igualdad es el teorema anterior, junto con \(S(f, P) \geq \upint_a^b f \,\dd x\) para toda \(P\). Para la segunda, \(s(f, P) = -S(-f, P)\) y \(\lowint_a^b f \,\dd x = -\upint_a^b (-f) \,\dd x\).
\end{proof}

\begin{definition}[Partición punteada]
  Una \emph{partición punteada} \(P^{*}\) de \([a, b]\) es una partición \(P = \{ t_0, \dots, t_n \}\) junto con una elección de puntos \(\alpha_i \in [t_{i-1}, t_i]\) para \(i \in \{1, \dots, n\}\). Su norma es \(\norm{P^{*}} = \norm{P}\).
\end{definition}

\begin{definition}[Suma de Riemann]
  Si \(f\colon [a, b] \to \R\) es acotada y \(P^{*}\) es una partición punteada, la \emph{suma de Riemann} de \(f\) asociada a \(P^{*}\) es
  \[
    R(f, P^{*}) = \sum_{i=1}^{n} f(\alpha_i)(t_i - t_{i-1}).
  \]
  Decimos que \(I = \lim_{\norm{P} \to 0} R(f, P^{*})\) si para todo \(\e > 0\) existe \(\delta > 0\) tal que \(|R(f, P^{*}) - I| < \e\) para toda partición punteada \(P^{*}\) con \(\norm{P^{*}} < \delta\).
\end{definition}

\begin{theorem}
  Sea \(f\colon [a, b] \to \R\) acotada. Existe \(I = \lim_{\norm{P} \to 0} R(f, P^{*})\) si y sólo si \(f\) es integrable, y en ese caso \(I = \int_a^b f(x) \,\dd x\).
\end{theorem}

\begin{proof}
  (\(\Leftarrow\)) Si \(f\) es integrable, por el corolario anterior \(\lim_{\norm{P} \to 0} s(f, P) = \lim_{\norm{P} \to 0} S(f, P) = \int_a^b f \,\dd x\). Como \(s(f, P) \leq R(f, P^{*}) \leq S(f, P)\) para toda partición punteada, \(\lim_{\norm{P} \to 0} R(f, P^{*}) = \int_a^b f \,\dd x\).

  (\(\Rightarrow\)) Supongamos que el límite \(I\) existe y sea \(\e > 0\). Existe \(\delta > 0\) tal que \(|R(f, P^{*}) - I| < \frac{\e}{4}\) para toda partición punteada con \(\norm{P^{*}} < \delta\). Fijemos una partición \(P = \{ t_0, \dots, t_n \}\) con \(\norm{P} < \delta\). Elegimos \(\alpha_i \in [t_{i-1}, t_i]\) con \(f(\alpha_i) < m_i + \frac{\e}{4n(t_i - t_{i-1})}\); entonces
  \[
    R(f, P^{*}) < \sum_{i=1}^{n} m_i (t_i - t_{i-1}) + \frac{\e}{4} = s(f, P) + \frac{\e}{4}.
  \]
  Análogamente, eligiendo otros puntos \(\tilde{\alpha}_i\) con \(f(\tilde{\alpha}_i) > M_i - \frac{\e}{4n(t_i - t_{i-1})}\), obtenemos una partición punteada \(\tilde{P}^{*}\) con \(R(f, \tilde{P}^{*}) > S(f, P) - \frac{\e}{4}\). Luego
  \[
    I - \frac{\e}{2} < R(f, P^{*}) - \frac{\e}{4} < s(f, P) \leq S(f, P) < R(f, \tilde{P}^{*}) + \frac{\e}{4} < I + \frac{\e}{2},
  \]
  de donde \(S(f, P) - s(f, P) < \e\). Por el teorema de caracterización, \(f\) es integrable. Además \(s(f, P) \leq \int_a^b f \,\dd x \leq S(f, P)\) y \(I\) pertenece al mismo intervalo \([s(f, P), S(f, P)]\), de longitud menor que \(\e\). Como \(\e\) es arbitrario, \(I = \int_a^b f(x) \,\dd x\).
\end{proof}

\begin{eg}
  Sea \(f\colon [1, 2] \to \R\), \(f(x) = \frac{1}{x}\). Por la regla de Barrow, \(\int_1^2 \frac{1}{x} \,\dd x = \ln 2 - \ln 1 = \ln 2\). Para cada \(n \in \N\) sea \(P_n\) la partición \(t_i = \frac{n + i}{n}\), con \(i \in \{0, \dots, n\}\), y elijamos \(\alpha_i = t_i\) en cada \([t_{i-1}, t_i]\). Entonces \(t_i - t_{i-1} = \frac{1}{n}\), \(f(\alpha_i) = \frac{n}{n + i}\) y
  \[
    R(f, P_n^{*}) = \sum_{i=1}^{n} \frac{n}{n + i} \cdot \frac{1}{n} = \sum_{i=1}^{n} \frac{1}{n + i}.
  \]
  Como \(\norm{P_n} = \frac{1}{n} \to 0\) y \(f\) es integrable,
  \[
    \lim_{n \to \infty} \left( \frac{1}{n + 1} + \cdots + \frac{1}{2n} \right) = \ln 2.
  \]
\end{eg}

\section{Integral de Riemann en varias dimensiones}

Un \emph{intervalo \(n\)-dimensional} es un conjunto \(I = I_1 \times \cdots \times I_n \subset \R^n\), donde \(I_1, \dots, I_n\) son intervalos de \(\R\). Se admite el caso degenerado en que algunos \(I_k\) se reducen a un punto. Si todos los \(I_k\) son abiertos (respectivamente cerrados, acotados), entonces \(I\) es abierto (respectivamente cerrado, acotado) en \(\R^n\). Notamos \(|I| = |I_1| \cdots |I_n|\), donde \(|I_k|\) es la longitud de \(I_k\).

Sea \(I = [a_1, b_1] \times \cdots \times [a_n, b_n]\) un intervalo compacto. Si \(P_k\) es una partición de \([a_k, b_k]\) para cada \(k\), entonces \(P = P_1 \times \cdots \times P_n\) es una \emph{partición} de \(I\), que determina subintervalos compactos \(Q_1, \dots, Q_N\) (los productos de subintervalos de las \(P_k\)), cuyos interiores son disjuntos dos a dos y cuya unión es \(I\). Decimos que una partición \(P' = P'_1 \times \cdots \times P'_n\) es \emph{más fina} que \(P\) si \(P_k \subset P'_k\) para todo \(k\).

\begin{definition}[Suma de Riemann e integrabilidad]
  Sea \(f\colon I \to \R\) acotada, \(P\) una partición de \(I\) con subintervalos \(Q_1, \dots, Q_N\) y \(\alpha_j \in Q_j\) para cada \(j\). La \emph{suma de Riemann} asociada es \(R(f, P^{*}) = \sum_{j=1}^{N} f(\alpha_j) |Q_j|\). Decimos que \(f\) es \emph{integrable Riemann} en \(I\) si existe \(A \in \R\) tal que para todo \(\e > 0\) existe una partición \(P_{\e}\) de \(I\) con \(|R(f, P^{*}) - A| < \e\) para toda partición \(P\) más fina que \(P_{\e}\) y toda elección de puntos \(\alpha_j \in Q_j\). En ese caso escribimos
  \[
    A = \int_I f(x) \,\dd x = \int \cdots \int_I f(x_1, \dots, x_n) \,\dd x_1 \cdots \dd x_n.
  \]
\end{definition}

\begin{definition}[Sumas e integrales de Darboux]
  Sea \(I \subset \R^n\) un intervalo compacto, \(f\colon I \to \R\) acotada y \(P\) una partición de \(I\) con subintervalos \(Q_1, \dots, Q_N\). Sean \(m_j = \inf f(Q_j)\) y \(M_j = \sup f(Q_j)\). Las sumas inferior y superior son \(s(f, P) = \sum_{j=1}^{N} m_j |Q_j|\) y \(S(f, P) = \sum_{j=1}^{N} M_j |Q_j|\), y las integrales inferior y superior son
  \[
    \lowint_I f(x) \,\dd x = \sup_P s(f, P), \qquad \upint_I f(x) \,\dd x = \inf_P S(f, P).
  \]
  Estas cantidades mantienen las propiedades vistas en \(\R\).
\end{definition}

\begin{theorem}
  Sea \(I \subset \R^n\) un intervalo compacto y \(f\colon I \to \R\) acotada. Son equivalentes:
  \begin{enumerate}[label=(\arabic*)]
    \item \(f\) es integrable en \(I\).
    \item Para todo \(\e > 0\) existe una partición \(P_{\e}\) de \(I\) tal que \(S(f, P) - s(f, P) < \e\) para toda \(P\) más fina que \(P_{\e}\).
    \item \(\lowint_I f(x) \,\dd x = \upint_I f(x) \,\dd x\).
  \end{enumerate}
\end{theorem}

\begin{proof}
  La demostración es análoga a la del caso de una variable y se deja como ejercicio.
\end{proof}

\begin{prop}
  Si \(I \subset \R^n\) es un intervalo compacto y \(f\colon I \to \R\) es continua, entonces \(f\) es integrable en \(I\).
\end{prop}

\begin{proof}
  Podemos suponer \(|I| > 0\). Como \(I\) es compacto, \(f\) es acotada y uniformemente continua. Dado \(\e > 0\) existe \(\delta > 0\) tal que \(\norm{x - y} < \delta\) implica \(|f(x) - f(y)| < \frac{\e}{2|I|}\). Sea \(P\) una partición cuyos subintervalos tienen diámetro menor que \(\delta\) (basta subdividir cada lado en partes pequeñas). Entonces \(M_j - m_j \leq \frac{\e}{2|I|}\) para todo \(j\) y
  \[
    S(f, P) - s(f, P) = \sum_{j} (M_j - m_j) |Q_j| \leq \frac{\e}{2|I|} \sum_j |Q_j| = \frac{\e}{2} < \e.
  \]
  Como \(s(f, P) \leq \lowint_I f \leq \upint_I f \leq S(f, P)\), resulta \(0 \leq \upint_I f - \lowint_I f < \e\) para todo \(\e > 0\), es decir, \(f\) es integrable.
\end{proof}

\begin{theorem}
  Sea \(Q = [a, b] \times [c, d] \subset \R^2\) y \(f\colon Q \to \R\) acotada. Entonces
  \[
    \lowint_Q f \leq \lowint_a^b \left( \upint_c^d f(x, y) \,\dd y \right) \dd x \leq \upint_a^b \left( \upint_c^d f(x, y) \,\dd y \right) \dd x \leq \upint_Q f.
  \]
\end{theorem}

\begin{proof}
  Sea \(F(x) = \upint_c^d f(x, y) \,\dd y\) para \(x \in [a, b]\), que está acotada por \(M(d - c)\) si \(|f| \leq M\). Sean \(P_1 = \{ x_0, \dots, x_p \}\) y \(P_2 = \{ y_0, \dots, y_q \}\) particiones de \([a, b]\) y \([c, d]\), y \(P = P_1 \times P_2\), con subrectángulos \(Q_{ij} = [x_{i-1}, x_i] \times [y_{j-1}, y_j]\). Sean \(m_{ij} = \inf f(Q_{ij})\) y \(M_{ij} = \sup f(Q_{ij})\).

  Por la aditividad de la integral superior, \(F(x) = \sum_{j=1}^{q} \upint_{y_{j-1}}^{y_j} f(x, y) \,\dd y\). Si \(x \in [x_{i-1}, x_i]\), entonces \(m_{ij} \leq f(x, y) \leq M_{ij}\) para \(y \in [y_{j-1}, y_j]\), y por lo tanto
  \[
    \sum_{j=1}^{q} m_{ij} (y_j - y_{j-1}) \leq F(x) \leq \sum_{j=1}^{q} M_{ij} (y_j - y_{j-1}).
  \]
  Tomando ínfimo y supremo de \(F\) en \([x_{i-1}, x_i]\) y multiplicando por \(x_i - x_{i-1}\), y sumando en \(i\), obtenemos
  \[
    s(f, P) \leq s(F, P_1) \qquad \text{y} \qquad S(F, P_1) \leq S(f, P).
  \]
  Entonces \(s(f, P) \leq s(F, P_1) \leq \lowint_a^b F \,\dd x \leq \upint_a^b F \,\dd x \leq S(F, P_1) \leq S(f, P)\). Como esto vale para toda partición \(P = P_1 \times P_2\) de \(Q\), tomando supremo sobre \(P\) en el miembro izquierdo e ínfimo en el derecho se obtiene el resultado.
\end{proof}

\begin{note}
  Con la misma demostración se obtiene lo siguiente.
  \begin{enumerate}[label=(\arabic*)]
    \item Vale el teorema reemplazando \(\upint_c^d\) por \(\lowint_c^d\):
          \[
            \lowint_Q f \leq \lowint_a^b \bigl( \lowint_c^d f \,\dd y \bigr) \dd x \leq \upint_a^b \bigl( \lowint_c^d f \,\dd y \bigr) \dd x \leq \upint_Q f.
          \]
    \item Vale intercambiando los papeles de \(x\) e \(y\):
          \[
            \lowint_Q f \leq \lowint_c^d \bigl( \upint_a^b f \,\dd x \bigr) \dd y \leq \upint_c^d \bigl( \upint_a^b f \,\dd x \bigr) \dd y \leq \upint_Q f,
          \]
          y también con \(\lowint_a^b\) en el interior.
    \item Si \(f\) es integrable en \(Q\), los extremos de todas estas cadenas de desigualdades coinciden con \(\iint_Q f\), y por lo tanto todas las integrales iteradas (con cualquier combinación de integrales inferiores o superiores) son iguales a \(\iint_Q f\).
  \end{enumerate}
\end{note}

\section{Teorema de Fubini}

\begin{corollary}[Teorema de Fubini]
  Si \(f\colon Q \to \R\) es continua en \(Q = [a, b] \times [c, d]\), entonces
  \[
    \iint_Q f = \int_a^b \left( \int_c^d f(x, y) \,\dd y \right) \dd x = \int_c^d \left( \int_a^b f(x, y) \,\dd x \right) \dd y.
  \]
\end{corollary}

\begin{proof}
  Por la proposición anterior \(f\) es integrable en \(Q\). Para cada \(x\), la función \(y \mapsto f(x, y)\) es continua, y por lo tanto integrable, de modo que \(\upint_c^d f(x, y) \,\dd y = \int_c^d f(x, y) \,\dd y\). Por la nota anterior (3), \(\lowint_a^b F \,\dd x = \upint_a^b F \,\dd x = \iint_Q f\), donde \(F(x) = \int_c^d f(x, y) \,\dd y\). Es decir, \(F\) es integrable y \(\iint_Q f = \int_a^b F(x) \,\dd x\). La otra igualdad se prueba intercambiando los papeles de \(x\) e \(y\).
\end{proof}

\section{Contenido de Jordan}

\begin{definition}[Contenidos interior y exterior]
  Sea \(S\) un subconjunto de un intervalo compacto \(I \subset \R^n\). Para cada partición \(P\) de \(I\) sea \(\underline{J}(S, P)\) la suma de las medidas de los subintervalos de \(P\) contenidos en el interior de \(S\), y \(\overline{J}(S, P)\) la suma de las medidas de los subintervalos de \(P\) que intersecan a \(\cl{S}\). El \emph{contenido interior} y el \emph{contenido exterior} de \(S\) son
  \[
    \underline{c}(S) = \sup_P \underline{J}(S, P), \qquad \overline{c}(S) = \inf_P \overline{J}(S, P).
  \]
\end{definition}

\begin{definition}[Contenido de Jordan]
  Si \(\underline{c}(S) = \overline{c}(S)\), este valor común se denota \(c(S)\) y se llama \emph{contenido de Jordan} de \(S\). En ese caso decimos que \(S\) es \emph{medible Jordan}.
\end{definition}

\begin{note}
  \begin{enumerate}[label=(\arabic*)]
    \item Se puede ver que \(\underline{c}(S)\) y \(\overline{c}(S)\) dependen únicamente de \(S\) y no del intervalo \(I\) elegido.
    \item \(0 \leq \underline{c}(S) \leq \overline{c}(S)\).
    \item Si \(c(S) = 0\), entonces \(\overline{c}(S) = 0\) y, por definición de ínfimo, para todo \(\e > 0\) el conjunto \(S\) puede cubrirse con finitos intervalos cuyas medidas suman menos que \(\e\).
  \end{enumerate}
\end{note}
